\documentclass[11pt]{article}
\usepackage[utf8]{inputenc}
\usepackage[margin=1in]{geometry}
\usepackage{natbib}
\usepackage{booktabs}
\usepackage{amsmath}
\usepackage{graphicx}
\usepackage{hyperref}
\usepackage{xcolor}

\title{Thiolutin Is a Direct Inhibitor of RNA Polymerase~II}

\author{
  Author One$^{1}$, Author Two$^{1}$, Author Three$^{2}$, Author Four$^{1,*}$ \\
  $^{1}$Department of Biochemistry, University A \\
  $^{2}$Department of Molecular Biology, University B \\
  $^{*}$Corresponding author
}

\date{}

\begin{document}
\maketitle

\begin{abstract}
Thiolutin, a natural product dithiolopyrrolone antibiotic, has been widely used as a transcription inhibitor, yet its mechanism of action has been debated for decades. Previous studies attributed its in vivo effects on transcription to indirect mechanisms involving zinc chelation and oxidative stress. Here we combine chemicogenetic screening in \textit{Saccharomyces cerevisiae} with in vitro transcription assays using purified yeast RNA Polymerase~II to demonstrate that thiolutin is a direct Pol~II inhibitor. A genome-wide deletion screen reveals that thioredoxin mutants (\textit{trx1$\Delta$}, \textit{trx2$\Delta$}) and metal transporter mutants (\textit{pmr1$\Delta$}, \textit{ctr1$\Delta$}) are hypersensitive to thiolutin, implicating both redox buffering and multi-metal interactions. In vivo, thiolutin causes strong thioredoxin oxidation (65--72\% oxidized vs.\ 8--12\% in controls). In vitro, thiolutin directly inhibits Pol~II transcription initiation with an IC$_{50}$ of $\sim$12~$\mu$M, but only in the presence of Mn$^{2+}$; removal of Mn$^{2+}$ or addition of excess DTT (5~mM) abrogates inhibition almost completely (88--92\% residual activity). The structurally related compound holomycin shows comparable direct inhibition (22\% residual transcription at 25~$\mu$M). These results reconcile conflicting literature by demonstrating that thiolutin possesses both direct Pol~II inhibitory activity and indirect effects through metal chelation and thiol oxidation, with the balance determined by intracellular redox and metal conditions.
\end{abstract}

\section{Introduction}

Thiolutin is a member of the dithiolopyrrolone family of natural antibiotics, first isolated from \textit{Streptomyces albus} \citep{celmer1952}. It has been widely used as a tool compound for rapid transcription shutoff in yeast and other organisms \citep{jimenez1973,tipper1973}. Despite decades of use, the molecular mechanism underlying thiolutin's transcription-inhibitory activity has remained controversial.

Early studies demonstrated that thiolutin inhibits RNA synthesis broadly, affecting RNA Polymerases I, II, and III in vitro \citep{tipper1973,khachatourians1975}. However, subsequent work by \citet{lauinger2017} argued that thiolutin's in vivo effects on transcription are indirect, mediated primarily through Zn$^{2+}$ chelation. The dithiolopyrrolone ring system contains sulfur atoms capable of coordinating metal ions, and thiolutin was shown to deplete intracellular zinc, disrupting zinc-dependent transcription factors \citep{lauinger2017}. This interpretation questioned the utility of thiolutin as a direct transcription inhibitor.

However, the zinc chelation model does not account for all observations. Thiolutin and the related compound holomycin are known to undergo redox cycling and to generate reactive oxygen species \citep{oliva2001,li2007}. The contribution of oxidative stress to thiolutin's biological effects has not been systematically dissected. Moreover, whether thiolutin interacts functionally with metals beyond zinc---particularly manganese, which is essential for certain polymerase activities---has not been examined.

Here we use a combination of chemicogenetic screening in \textit{S.~cerevisiae} and reconstituted in vitro transcription assays to revisit the mechanism of thiolutin. We identify genetic determinants of thiolutin sensitivity that implicate both redox homeostasis and multi-metal interactions, and demonstrate that thiolutin directly inhibits RNA Polymerase~II transcription initiation in a manganese-dependent and redox-sensitive manner.

\section{Related Work}

\subsection{Dithiolopyrrolone Antibiotics}

The dithiolopyrrolone family includes thiolutin, holomycin, and related compounds produced by various \textit{Streptomyces} species \citep{celmer1952,kenig1979}. These compounds exhibit broad-spectrum antimicrobial activity and have been shown to affect multiple cellular processes including transcription, translation, and cell wall biosynthesis \citep{oliva2001}. The core bicyclic ring system containing a disulfide bridge is essential for biological activity and undergoes redox cycling between oxidized and reduced forms \citep{li2007}.

\subsection{Thiolutin as a Transcription Inhibitor}

Thiolutin's use as a transcription shutoff agent in yeast dates to the 1970s \citep{jimenez1973,tipper1973}. The compound was shown to inhibit all three nuclear RNA polymerases in vitro \citep{khachatourians1975}. More recently, \citet{lauinger2017} challenged the direct inhibition model, proposing that zinc chelation accounts for the in vivo transcription defect. However, the in vitro conditions used in these studies varied considerably, particularly regarding metal ion composition and reducing agents, potentially explaining discrepant results.

\subsection{Chemicogenomic Approaches in Yeast}

Genome-wide deletion library screens have proven powerful for elucidating drug mechanisms of action \citep{giaever2007,hillenmeyer2008}. Chemicogenetic profiling---identifying mutants with altered drug sensitivity---reveals functional connections between drug targets and cellular pathways \citep{costanzo2019}. This approach has been applied to numerous compounds but not systematically to thiolutin.

\section{Methods}

\subsection{Chemicogenetic Screen}

The \textit{S.~cerevisiae} homozygous deletion collection was screened for growth on YPD agar supplemented with 5~$\mu$g/mL thiolutin or DMSO vehicle control. Colony sizes were measured after 48~h incubation at 30$^{\circ}$C and growth ratios (thiolutin/DMSO) were calculated for each mutant. Hits were defined as mutants with growth ratios $<$0.3 (sensitive) or $>$0.85 (resistant).

\subsection{Thioredoxin Oxidation Assay}

Wild-type yeast were treated with 3~$\mu$g/mL thiolutin or DMSO for 15~min. Cell extracts were prepared under denaturing conditions and thioredoxin (Trx1, Trx2) redox state was assessed by 4-acetamido-4$'$-maleimidylstilbene-2,2$'$-disulfonic acid (AMS) gel shift assay \citep{berndt2008}. H$_2$O$_2$ (1~mM, 15~min) served as positive control.

\subsection{Metal Supplementation Assays}

Growth inhibition by thiolutin (3~$\mu$g/mL) was measured in liquid YPD supplemented with 100~$\mu$M of individual metal ions (ZnCl$_2$, MnCl$_2$, CuSO$_4$, FeSO$_4$, MgCl$_2$). Growth was measured by OD$_{600}$ after 16~h and expressed as percentage of untreated control. Statistical significance was assessed by Student's $t$-test.

\subsection{In Vitro Transcription Assay}

Yeast RNA Polymerase~II was purified from \textit{S.~cerevisiae} as described \citep{cramer2000}. Transcription initiation was reconstituted on a model promoter template containing a HIS4 promoter fused to a G-less cassette \citep{lorch1997}. Reactions contained purified Pol~II, general transcription factors (TFIIB, TFIID, TFIIE, TFIIF, TFIIH), 1~mM MnCl$_2$ (where indicated), 0.5~mM DTT (standard) or 5~mM DTT (excess), and radiolabeled [$\alpha$-$^{32}$P]CTP. Runoff transcripts were resolved by denaturing PAGE and quantified by phosphorimaging. Thiolutin or holomycin was added at specified concentrations prior to NTP addition. IC$_{50}$ values were determined by nonlinear regression fitting.

\section{Results}

\subsection{Chemicogenetic Screen Identifies Redox and Metal Homeostasis Genes}

We screened the yeast deletion library for altered sensitivity to thiolutin at 5~$\mu$g/mL (Table~\ref{tab:screen}). The most sensitive mutants included both thioredoxin genes (\textit{trx1$\Delta$}, ratio 0.12; \textit{trx2$\Delta$}, ratio 0.15) and the superoxide dismutase gene \textit{sod1$\Delta$} (ratio 0.22), strongly implicating redox buffering as a critical determinant of thiolutin tolerance. Notably, the Mn$^{2+}$ transporter mutant \textit{pmr1$\Delta$} (ratio 0.18) and the Cu$^{2+}$ transporter mutant \textit{ctr1$\Delta$} (ratio 0.25) were also hypersensitive, indicating that thiolutin interacts functionally with multiple metals beyond zinc. Among resistant mutants, the zinc transporter \textit{zrt1$\Delta$} (ratio 1.05) showed the strongest resistance, consistent with the established zinc chelation component of thiolutin activity.

\begin{table}[ht]
\centering
\caption{Top chemicogenetic hits for thiolutin sensitivity and resistance.}
\label{tab:screen}
\begin{tabular}{llrl}
\toprule
\textbf{Gene} & \textbf{Function} & \textbf{Growth Ratio} & \textbf{Interpretation} \\
\midrule
\multicolumn{4}{l}{\textit{Sensitive mutants}} \\
\textit{TRX1} & Thioredoxin & 0.12 & Redox buffering \\
\textit{TRX2} & Thioredoxin & 0.15 & Redox buffering \\
\textit{PMR1} & Mn$^{2+}$ transporter & 0.18 & Mn involvement \\
\textit{SOD1} & Superoxide dismutase & 0.22 & Oxidative stress \\
\textit{CTR1} & Cu$^{2+}$ transporter & 0.25 & Cu involvement \\
\midrule
\multicolumn{4}{l}{\textit{Resistant mutants}} \\
\textit{ZRT1} & Zn$^{2+}$ transporter & 1.05 & Zn chelation \\
\textit{FET3} & Fe/Cu oxidase & 0.95 & Metal homeostasis \\
\bottomrule
\end{tabular}
\end{table}

\subsection{Thiolutin Causes Thioredoxin Oxidation In Vivo}

To directly test whether thiolutin perturbs cellular redox state, we measured thioredoxin oxidation by AMS gel shift (Table~\ref{tab:trx}). Treatment with 3~$\mu$g/mL thiolutin for 15~min increased the fraction of oxidized Trx1 from 8\% to 65\% and oxidized Trx2 from 12\% to 72\%. These levels approach those observed with the H$_2$O$_2$ positive control (78\% and 81\%, respectively). This confirms that thiolutin is a potent thiol oxidant in vivo and explains the extreme sensitivity of thioredoxin mutants.

\begin{table}[ht]
\centering
\caption{Thioredoxin oxidation upon thiolutin treatment.}
\label{tab:trx}
\begin{tabular}{lrr}
\toprule
\textbf{Condition} & \textbf{\% Oxidized Trx1} & \textbf{\% Oxidized Trx2} \\
\midrule
DMSO control & 8 & 12 \\
Thiolutin (3~$\mu$g/mL) & 65 & 72 \\
H$_2$O$_2$ control & 78 & 81 \\
\bottomrule
\end{tabular}
\end{table}

\subsection{Thiolutin Interacts Functionally with Multiple Metals}

Metal supplementation experiments revealed that both Zn$^{2+}$ and Mn$^{2+}$ significantly rescued thiolutin-induced growth inhibition (Table~\ref{tab:metals}). Zn$^{2+}$ restored growth from 35\% to 78\% ($p < 0.001$), and Mn$^{2+}$ restored growth to 70\% ($p < 0.001$). In contrast, Cu$^{2+}$ worsened growth inhibition (28\%), while Fe$^{2+}$ and Mg$^{2+}$ had no significant effect. These results demonstrate that thiolutin's biological activity involves interactions with at least two essential metals, not only zinc as previously proposed.

\begin{table}[ht]
\centering
\caption{Metal supplementation rescue of thiolutin growth inhibition.}
\label{tab:metals}
\begin{tabular}{lrl}
\toprule
\textbf{Metal (100~$\mu$M)} & \textbf{Growth (\% of untreated)} & \textbf{$p$-value} \\
\midrule
None (thiolutin only) & 35 & -- \\
Zn$^{2+}$ & 78 & $<$0.001 \\
Mn$^{2+}$ & 70 & $<$0.001 \\
Cu$^{2+}$ & 28 & n.s.\ (worsened) \\
Fe$^{2+}$ & 40 & n.s. \\
Mg$^{2+}$ & 36 & n.s. \\
\bottomrule
\end{tabular}
\end{table}

\subsection{Thiolutin Directly Inhibits RNA Pol~II Transcription In Vitro}

Using purified yeast Pol~II in a reconstituted transcription initiation assay, we tested whether thiolutin directly inhibits transcription (Table~\ref{tab:invitro}). In the presence of 1~mM Mn$^{2+}$ and standard reducing conditions (0.5~mM DTT), thiolutin inhibited Pol~II transcription in a dose-dependent manner: 72\% residual activity at 5~$\mu$M, 45\% at 10~$\mu$M, 18\% at 25~$\mu$M, and 8\% at 50~$\mu$M, yielding an IC$_{50}$ of approximately 12~$\mu$M.

Critically, this inhibition was strictly dependent on both Mn$^{2+}$ and appropriate redox conditions. In the absence of Mn$^{2+}$, 25~$\mu$M thiolutin caused only minimal inhibition (92\% residual transcription). Similarly, excess DTT (5~mM) almost completely abrogated inhibition (88\% residual activity at 25~$\mu$M thiolutin). The structurally related dithiolopyrrolone holomycin showed comparable inhibitory activity (22\% residual transcription at 25~$\mu$M with Mn$^{2+}$), indicating that direct Pol~II inhibition is a general property of this compound class.

\begin{table}[ht]
\centering
\caption{In vitro Pol~II transcription inhibition by thiolutin under varying conditions.}
\label{tab:invitro}
\begin{tabular}{lrll}
\toprule
\textbf{Thiolutin ($\mu$M)} & \textbf{Rel.\ Transcription (\%)} & \textbf{Mn$^{2+}$} & \textbf{DTT} \\
\midrule
0 & 100 & 1~mM & 0.5~mM \\
5 & 72 & 1~mM & 0.5~mM \\
10 & 45 & 1~mM & 0.5~mM \\
25 & 18 & 1~mM & 0.5~mM \\
50 & 8 & 1~mM & 0.5~mM \\
25 & 92 & None & 0.5~mM \\
25 & 88 & 1~mM & 5~mM \\
\bottomrule
\end{tabular}
\end{table}

\section{Discussion}

Our findings resolve a long-standing controversy regarding thiolutin's mechanism of action by demonstrating that it possesses genuine direct inhibitory activity against RNA Polymerase~II, while also confirming the indirect contributions of metal chelation and oxidative stress observed by others \citep{lauinger2017}.

The chemicogenetic screen provides an unbiased view of the cellular pathways affected by thiolutin. The extreme sensitivity of thioredoxin mutants, combined with the demonstration that thiolutin causes near-complete thioredoxin oxidation in vivo, establishes oxidative stress as a major component of thiolutin's cellular effects. This is consistent with the known redox-active disulfide bridge in the dithiolopyrrolone ring \citep{li2007,oliva2001}. The sensitivity of the Mn$^{2+}$ transporter mutant \textit{pmr1$\Delta$} extends the metal interaction model beyond zinc to include manganese, a finding with important mechanistic implications.

The in vitro transcription data provide definitive evidence for direct Pol~II inhibition. The IC$_{50}$ of $\sim$12~$\mu$M is within the physiologically relevant range given the intracellular concentrations achieved at the doses used in vivo (3--5~$\mu$g/mL) \citep{jimenez1973}. The strict requirement for Mn$^{2+}$ suggests that thiolutin may form a metal-coordinated complex that interacts with the Pol~II active site, which itself contains two catalytic metal ions \citep{cramer2001,gnatt2001}. The abrogation of inhibition by excess DTT indicates that the oxidized (disulfide-containing) form of thiolutin is the active inhibitory species, consistent with recent structural studies of dithiolopyrrolone redox chemistry \citep{li2007}.

The reconciliation of direct and indirect mechanisms has practical implications for the use of thiolutin as a transcription shutoff tool. Under standard laboratory conditions, both mechanisms likely contribute to the observed rapid transcription shutdown. The relative contribution of each mechanism will depend on the intracellular redox state and metal ion availability of the cell type under study.

Our work also establishes holomycin as a direct Pol~II inhibitor with comparable potency, expanding the potential pharmacological applications of the dithiolopyrrolone compound class \citep{chan2017}.

\section{Conclusion}

We demonstrate that thiolutin is a direct, Mn$^{2+}$-dependent, redox-sensitive inhibitor of RNA Polymerase~II transcription initiation, with an IC$_{50}$ of $\sim$12~$\mu$M. Chemicogenetic profiling reveals that thiolutin sensitivity is governed by redox homeostasis and multi-metal interactions. These findings reconcile decades of conflicting literature and reaffirm thiolutin's utility as a transcription inhibitor while defining the critical conditions---metal ions and redox environment---that determine its direct versus indirect modes of action.

\bibliographystyle{plainnat}
\bibliography{references}

\end{document}
